\begin{papercard}{brooks2017}
\cf{Stages} \stagetag{S} (privacy problem, privacy engineering, problematic data action), \stagetag{H} (risk = likelihood $\times$ impact of a problem for individuals; data actions and context). NISTIR 8062 (January 2017, 49 pages), an introductory federal report that proposes (i) a definition of privacy engineering, (ii) three privacy engineering objectives (predictability, manageability, disassociability), and (iii) a privacy risk model built around the \emph{problematic data action}. It is guidance, not a study: no data, no test of the model. Its value is a government-grade vocabulary for the step from a data practice to a risk to individuals.

\cf{Motivation} NIST security guidance addresses confidentiality, integrity and availability, but ``there is a far less developed understanding of how the system impacts the individuals' privacy and how to identify and address risks that extend beyond unauthorized access to PII'' \pc{brooks2017}{9}{there is a far less developed understanding of how the system impacts the individuals’ privacy and how to identify and address risks that extend beyond unauthorized access to PII}. Examples: smart meters (house occupancy and appliance use); the Veterans' genomic programme whose privacy impact assessment lists re-identification, participants who ``misunderstand or underestimate the extent to which they have consented to share their data'', stigma, discrimination and law-enforcement access \pc{brooks2017}{10}{Participants misunderstand or underestimate the extent to which they have consented to share their data}. Privacy concerns can arise at any stage of PII processing, ``collection, creation, analysis, use, storage, dissemination, disclosure, or disposal'' \pc{brooks2017}{10}{privacy concerns can arise at any point where PII is processed, including, but not limited to collection, creation, analysis, use, storage, dissemination, disclosure, or disposal}.

\cf{Security versus privacy: the byproduct risk model} In the security model the concern is unauthorized activity; in the privacy ``byproduct'' model ``the processing of PII is planned and permissible (i.e.\ authorized), but it creates implications for individuals' privacy'' \pc{brooks2017}{16}{In the byproduct risk model, the processing of PII is planned and permissible (i.e. authorized), but it creates implications for individuals’ privacy}. Security and privacy overlap (confidentiality of PII) but neither contains the other (Figure 2). Existing assessments tend to check compliance, not effect: ``assessments are conducted to determine whether a notice regarding privacy exists rather than to evaluate whether people are likely to read that notice and receive some privacy-protective benefit'' \pc{brooks2017}{22}{Assessments therefore, tend to be focused on compliance rather than the effectiveness of achieving a positive outcome for privacy}.

\cf{Privacy problems} The report uses the neutral term ``privacy problem'' for an adverse experience resulting from PII processing (alternatives: harm, violation, intrusion, invasion) \pc{brooks2017}{17}{The term “privacy problem” to express an adverse experience resulting from PII processing is not the only term that may be used}. Problems from unauthorised access are well recognised (embarrassment, identity theft, stalking); those from authorised processing are less visible: discriminatory and stigmatising effects of broad collection, inability to correct inaccurate data, and ``a loss of trust that results in reluctance to adopt certain products and services that could be beneficial'' \pc{brooks2017}{17}{Concerns about privacy and systems can cause a loss of trust that results in reluctance to adopt certain products and services that could be beneficial}. Figure 3 (derived from Solove's taxonomy) shows a range of problems, offered as illustration only \pc{brooks2017}{18}{Figure 3 is for illustrative purposes only. It does not purport to capture the full scope of problems people could experience from the processing of their information}.

\cf{Definition: privacy engineering} ``A specialty discipline of systems engineering focused on achieving freedom from conditions that can create problems for individuals with unacceptable consequences that arise from the system as it processes PII'' \pc{brooks2017}{6}{privacy engineering means a specialty discipline of systems engineering focused on achieving freedom from conditions that can create problems for individuals with unacceptable consequences that arise from the system as it processes PII}. It gives an outcome-based frame of reference; the FIPPs are ``value statements rather than recipes'' that are hard to operationalise \pc{brooks2017}{19}{They [the FIPPs] are value statements rather than recipes, however}. The phrase ``unacceptable consequences'' acknowledges that no system is free of privacy risk; risk management (NIST RMF, SP 800-30) is the process for deciding what is unacceptable.

\cf{Privacy engineering objectives (3.1)} Analogous to the CIA triad, three objectives characterise system capabilities. \emph{Predictability}: ``enabling of reliable assumptions by individuals, owners and operators about PII and its processing''; assess a notice by whether users read, understood and responded as anticipated \pc{brooks2017}{26}{an assessment of a notice as a means of enabling reliable assumptions might focus on whether users have read and understood the notice, or even whether they responded as anticipated}; it does not require users to know every technical detail but that they are not surprised \pc{brooks2017}{26}{Enabling reliable assumptions does not require that each stakeholder knows all the technical details about how a system processes PII}. \emph{Manageability}: granular administration of PII (alteration, deletion, selective disclosure); supports access and amendment, minimisation, quality and individual participation \pc{brooks2017}{27}{Manageability is an important system property enabling several of the FIPPs}; it is a capability, not a policy that individuals must have control. \emph{Disassociability}: processing without association to individuals or devices beyond operational need; ``unlike confidentiality ... disassociability recognizes that privacy risks can result from exposures even within an authorized perimeter'' \pc{brooks2017}{28}{Unlike confidentiality, which is focused on preventing unauthorized access to information, disassociability recognizes that privacy risks can result from exposures even within an authorized perimeter}. Figure 8 aligns the OMB Circular A-130 FIPPs (access and amendment, accountability, authority, minimisation, quality and integrity, individual participation, purpose specification and use limitation, security, transparency) with these objectives.

\cf{Privacy risk model (3.2)} In security, risk factors are threats and vulnerabilities; but ``the terms threat and vulnerability fail to capture the essence of many privacy problems for individuals'' because an authorised smart meter is not a ``threat''. The model therefore centres on the operation: ``Problematic data action means a data action that causes an adverse effect, or problem, for individuals'' \pc{brooks2017}{29}{Problematic data action means a data action that causes an adverse effect, or problem, for individuals}. \emph{Likelihood} is ``the probability that a data action will become problematic for a representative or typical individual whose PII is processed by the system'', informed by demographics, similar scenarios, focus groups and surveys \pc{brooks2017}{30}{Likelihood in a privacy risk model is the probability that a data action will become problematic for a representative or typical individual whose PII is processed by the system}. \emph{Impact} is ``the magnitude of cost or harm'', difficult to quantify because only individuals can experience the problem and experiences diverge; proxies are legal compliance, mission failure (reluctance to use), reputational and internal-culture costs \pc{brooks2017}{30}{The fact that only individuals—not agencies—can directly experience a privacy problem is especially challenging for assessing impact}. Characteristics: \emph{data actions} (collection, retention, logging, analysis, generation, transformation or merging, disclosure, transfer, disposal) around which an assessment should be organised \pc{brooks2017}{31}{a privacy risk assessment should be oriented around the identification of a system’s discrete data actions, and subsequent determination of which of these data actions could be problematic}; \emph{PII} (OMB: information that distinguishes or traces identity, alone or linked/linkable); \emph{context} (with Nissenbaum): an email address linked to a person with a health condition and disclosed outside health care ``may lead to problems such as embarrassment or discrimination'' \pc{brooks2017}{31}{an email address linked to an individual with a health condition and disclosed outside of the health care context may lead to problems such as embarrassment or discrimination}.

\cf{Roadmap and process} NIST plans guidance parallel to its security risk-management publications; its existing guidance is suited to harms from unauthorised access, less to authorised processing \pc{brooks2017}{32}{this guidance is appropriate for addressing risks to individuals arising from unauthorized access to their information}. Development process (Appendix A): public workshops (April 2014 about 240 attendees in person; later about 130 plus 500 webcast), a draft in May 2015 and over 40 comment submissions; the April workshop identified the communication gap between legal/policy and engineering teams \pc{brooks2017}{34}{There is a communication gap around privacy between the legal and policy, design and engineering, and product and project management teams}. Appendix E gives non-breach examples (cyber-physical systems and others), Appendix F the FIPPs.

\cf{Limits (card)} A federal introduction ``intended to foster further discussion''; the risk model is conceptual (no scales, no scoring method, no case application); US federal systems and Circular A-130 context; focus on organisational assessment, not on users' experience or evidence from apps; impact measured by proxies that favour the agency's costs; the objectives are not empirically validated; later NIST privacy guidance is not in the corpus and should be checked before citing this as current.

\cf{Use for this thesis} (1) The notion of a \emph{problematic data action} fits stage H: each finding is an app data action (collection, sharing, retention, inference, disclosure) described with its context, and the harm tree states what problem for which individual it can cause, rather than listing violations of principles. (2) The byproduct-risk idea supports the thesis stance that a practice can be authorised (consented, labelled) and still problematic; this matters for apps whose label or policy is formally compliant. (3) Predictability is a usable evaluation question for consent and notice screens: would a typical user be surprised by what happens to the data? and is the notice assessed by understanding, not mere presence? (4) Likelihood/impact language can structure the \emph{qualitative} severity discussion but the thesis should not compute risk scores, since the report offers no scoring method. (5) Cite together with Solove (taxonomy), Citron and Solove (harm types) and Nissenbaum (context).
\end{papercard}
